
\section{Background}

    A fundamental aspect of nature is that structures grow by applying built-in
    rules and adapting to circumstances rather than being constructed directly
    in a final state. It is very hard to get around this property, at least if
    we want to produce a large variety of organisms from the same base where
    the final configurations are different and adapted enough to appear
    convincing in their environment.



\subsection{Existing approaches}

    A number of system and approaches exists for modeling and evolving
    plants and other biological structures, most of them are based on
    Lindenmayer systems, described below. 



\subsection{Fractal properties}
    A common question about this project has been that of whether fractals are used.
    In order to answer that, an attempt is made to explain
    very briefly what a fractal actually is.\\
    Mathematically, the concept of a fractal dimension of a
    set is based upon the Housdorff-Besicovitch dimension,
    which is defined using the Hausdorff measure. It involves
    calculating the minimum number of a certain kind of sets
    required to cover the set of interest. There are
    heuristic ways to calculate this, amongst them the
    self-similarity dimension that makes use of the scaling
    property of a measure.\\
    A simpler and more popular way to define fractal is to
    say that it is "a geometric pattern repeated on every
    scale and so cannot be represented using classical
    geometry."\\
    Although somewhat dim, we could say that our trees hold
    some fractal properties, since they would appear when
    looked upon from a reasonable distance, to possess such
    repetetive patterns. These could be further extended by
    setting the the shapes of the branches to an intricate
    self-similar function.
 




\subsection{Lindenmayer System Introduction}

\label{lindenmayer}

    Aristid Lindenmayer proposed L-systems as a mathematical
    formalization of biological development in 1968. They have since
    been successfully used in computer graphics to generate realistic
    models of plants. 

    L-systems build on a very simple idea, successive rewriting of a
    string of symbols where the rules are applied in
    parallel\footnote{A similar system is Chomsky grammars, but the
    rules in these are applied sequentially, which does not reflect
    the parallel workings of nature as well as L-Systems.}.

    An L-system is described by an alphabet of symbols, a starting
    symbol and a number of rules. Originating from the starting symbol
    a string is built by applying matching rules iteratively.

    Consider the alphabet ${ab}$, starting symbol $a$ and rewrite
    rules ${a\rightarrow ab, b\rightarrow a}$. This successively produces the strings: 
    $a$, $ab$, $aba$, $abaab$, $abaababa$, \ldots 

    Now, of what use is strings of letters to model nature?
    The trick is to interpret the string as a geometrical description,
    one common way, but by no means the only one, is to use turtle graphics
    where symbols are interpreted either as lines to be
    drawn or commands to change orientation, simplified.

\begin{figure}[h]
\centering
\caption{Simple L-system }
\label{lsys-example}
\end{figure}
    
    Example \ref{lsys-example} is a D0L system, deterministic, context free
    (taking 0 neighbours in account), L-system. In context dependent
    systems (D$N$L) a rule may take a number of neighbouring symbols
    into account when deciding if it matches or not, however only a
    single symbol is rewritten.


\subsubsection{Parametric L-systems}

\label{param_l_system}

    In parametric L-Systems a number of parameters and
    conditions on these are also included in the rules, allowing for
    simpler descriptions of what would otherwise be quite
    incomprehensible (parameters can be simulated by extending the
    alphabet).

\begin{figure}[h]
\centering
\begin{tabular}{cc}
\multicolumn{2}{c}{
$ \begin{array}{lll}
\omega : & A(100,\omega_0) \\
p_1    : & A(s,w) : s \geq \mathit{min} \rightarrow & !(w)F(s) \\
 & & [+(\alpha_1)/(\varphi_1)A(s \cdot r_1, w \cdot q^e)] \\
 & & [+(\alpha_2)/(\varphi_2)A(s \cdot r_2, w \cdot (1-q)^e)]
\end{array} $ } \\
\\
\includegraphics[height=1in]{images/lsystem1} &
\includegraphics[height=1in]{images/lsystem2} \\
\end{tabular}
\caption{Parametric L-system (using notation from
\cite{lsys_theory_vis}) and possible visualizations, for two sets of
parameters}
\label{param-lsys-example}
\end{figure}

\subsubsection{Further extensions}

    Originating from this very simple idea, the L-systems have been
    extended and improved to support several interesting features such
    as signaling and bug-attacks. However, they tend to get
    complicated to grasp as they grow in complexity. One idea early on
    was simply to write a shell around this framework, to make it
    easier and more intuitive to use. It was however concluded that
    this approach might have caused unnecessary limitations. Instead
    it was agreed to begin at a higher level and perhaps later export
    to L-systems syntax.



